\section{Discussion and Threats to Validity}
\label{sec:discussion}

\subsection{Why Safety Obligations Help}

Safety obligations give static analysis and LLM reasoning different jobs.
Static analysis answers where a documented duty can arise and which program
object it concerns. The LLM answers a narrower semantic question: whether the
visible paths satisfy one of several explicit discharge conditions. This split
avoids requiring a complete handwritten transfer model for every API and
project helper, while preventing the model from choosing an arbitrary bug
property after reading the code.

The representation also makes reports comparable across surface forms. A
reference returned, transferred, or decremented uses different syntax but
resolves the same ownership duty. Conversely, a decrement can be correct for a
new reference and unsafe for a borrowed one. Encoding the required outcome and
its alternatives captures this distinction more directly than an isolated
pattern.

\subsection{Role of Static Pruning}

Static pattern-based pruning controls which obligation instances reach the LLM
rather than defining the safety semantics. A pruned instance is one where the
static evidence for a violation is absent; it may still contain a real defect
if that evidence is hidden inside a macro, a callee, or a generated wrapper.

\subsection{Scope and Limitations}

\tool is scoped to duties induced by the Python/C API. It does not attempt to
replace a general C/C++ memory-safety analyzer. The PyRefcon-only defects in
RQ2 provide concrete examples: raw allocation and null-dereference invariants
without a Python/C trigger never reach obligation binding. Extending the
method to project-local APIs would require another specification source and a
clear rule for composing local and CPython obligations.

The current context is intraprocedural source text. Function-like macros can
hide steal semantics, C++ destructors can hide release, and paired teardown
functions can hide ownership transfer. Of the 18 false positives, 2 are caused
by macro expansion, 2 by C++ RAII or destructor cleanup, and 2 by
interprocedural teardown helpers; the remaining 12 require richer
ownership-semantics knowledge to resolve. Preprocessing selected macros,
summarizing local helpers, and modeling common RAII wrappers could expose the
missing discharge evidence without broadening the property scope.

The structured checker remains an LLM-based judgment. It can omit a path,
accept infeasible control flow, or misunderstand an API. Temperature 0 reduces
sampling variation but does not provide determinism across hosted model
versions. The evidence-bearing schema supports audit, not formal soundness.
The high-confidence filter also trades report volume for review cost. The
current precision therefore describes the surfaced high-confidence set and
does not establish recall over all bugs or all model outputs.

\subsection{Threats to Validity}

\paragraph{Construct validity.}
The main metric is precision over high-confidence reports. This metric reflects
developer-facing audit cost but omits lower-confidence outputs and cannot
measure missed defects. PyCGuard and PyRefcon also use different report units
and filters. We state both units and use the comparison to study precision,
scope, and overlap rather than relative recall.

\paragraph{Internal validity.}
Manual labels can be wrong, particularly for rare failure paths and hidden
ownership conventions. The audit examines source context, API documentation,
and macros or callees when needed. Maintainer feedback and fixes provide
additional evidence for submitted reports, but feedback is not available for
every finding. The artifact includes report-level labels and
rationales so that independent reviewers can challenge individual decisions.

\paragraph{External validity.}
The subjects emphasize mature scientific, machine-learning, serialization,
networking, imaging, and system libraries. Results may not generalize to small
extensions, closed-source code, alternative Python implementations, or binding
frameworks that generate most native code. The obligations follow CPython 3.14
documentation. API evolution can add triggers or change conditions, requiring
the specification and index to be versioned.

\paragraph{Reproducibility.}
Hosted model behavior and legacy baseline toolchains can change. We report
exact model identifiers and decoding parameters and preserve prompts and raw
responses. The PyRefcon and CpyChecker environments are containerized, but
compilation still depends on project headers and generated files. The
anonymous artifact will include target manifests, containers, raw warnings,
deduplication records, and audit labels, together with an explanation for any
component that cannot be redistributed.
